\section{Structure of the Coupling}\label{sec:structure}

This section derives the one structural fact on which the receiver is built and the
consequences that each receiver block uses.

\begin{proposition}[Timing--propagation separation]\label{prop:sep}
For a piecewise-constant $K$-slot pattern with period $T_u$,
\begin{align}
b_q^{[h]}(f)&=\sum_{s=1}^{2}t_{q,s}[h]\,\Gamma_s(f),\nonumber\\
t_{q,s}[h]&=\frac{e^{-j\pi h/K}}{K}\,\mathrm{sinc}\!\Big(\frac hK\Big)\sum_{k:\,s_q[k]=s}e^{-j2\pi hk/K}.
\label{eq:tqs}
\end{align}
Hence, with $\bm a^{\rm d}_q=[a^{\rm d}_q(f_m)]_m$, $\bm a^{\rm s}_q=[a^{\rm s}_q(f_m)]_m$,
$\bm\Gamma_s=[\Gamma_s(f_m)]_m$ and Toeplitz matrices $[\bm T_{q,s}]_{\bar m m}=t_{q,s}[\bar m-m]$,
\begin{equation}
\bm H=\sum_{q=1}^{Q}\sum_{s=1}^{2}\bm T_{q,s}\circ\big(\bm a^{\rm d}_q(\bm\Gamma_s\odot\bm a^{\rm s}_q)^{\top}\big).
\label{eq:sep}
\end{equation}
\end{proposition}
\begin{proof}
During slot $k$ the element applies $\Gamma_{s_q[k]}(f)$, so its switched response is
$\sum_s\Gamma_s(f)\,\mathbb 1_s(t)$, where $\mathbb 1_s(t)$ indicates the slots in state $s$.
The $h$-th Fourier coefficient of the indicator of slot $k$ is
\[\tfrac1K e^{-j\pi h/K}\,\mathrm{sinc}(h/K)\,e^{-j2\pi hk/K};\]
summing over the slots in state $s$ gives $t_{q,s}[h]$. This is the slot sum of~\cite[eq.~(5)]{iudice2026} grouped by state.
Substituting into~\eqref{eq:io} and noting that $t_{q,s}$ depends on $\bar m-m$ only gives~\eqref{eq:sep}.
\end{proof}

The identity is elementary; its value lies in the following consequences.

\begin{corollary}[No element model is needed]\label{cor:nomodel}
The element response appears in~\eqref{eq:sep} only through $\bm\Gamma_s\odot\bm a^{\rm s}_q$,
a per-state \emph{input-side} response that also contains the unknown BS--RIS channel.
Estimating $\bm v_{q,s}:=\bm\Gamma_s\odot\bm a^{\rm s}_q$ as an unknown absorbs the dispersion
exactly; the receiver needs only $\bm T_{q,s}$, i.e. the switching timing.
\end{corollary}

\begin{corollary}[Fast products]\label{cor:fft}
$\big(\bm T\circ(\bm a\bm b^{\top})\big)\bm x=\bm a\odot\big(\bm T(\bm b\odot\bm x)\big)$, and
$\bm T\bm z$ is a linear convolution of length-$M$ sequences. Hence $\bm H\bm x$ and
$\bm H^{\mathsf H}\bm y$ cost $\mathcal O(QM\log M)$ with FFTs of length $\ge 3M-2$, and
$\bm H$ never needs to be formed. If all elements use the same pattern, $\bm T_{q,s}=\bm T_s$
for all $q$, and one batched FFT per state serves all elements.
\end{corollary}

\begin{corollary}[Harmonic tail]\label{cor:tail}
$|t_{q,s}[h]|\le\frac1K|\mathrm{sinc}(h/K)|\,|\{k:s_q[k]=s\}|$ decays as $1/|h|$, so the
coupling energy decays as $1/h^2$. For the ideal $\pm1$ element, $K=8$ and $D=0.75$, the
fraction of energy outside $|h|\le H$ is 9.7\%, 5.0\%, 3.4\%, 1.0\% and 0.31\% for
$H=4,8,12,40,128$ (numerical evaluation of~\eqref{eq:tqs}). A model truncated to
$|h|\le H$ therefore has a normalized error floor of at least $10\log_{10}$ of that fraction:
$-14.7$\,dB at $H=12$, which matches the $-14$ to $-16$\,dB floor measured for a truncated
estimator (Section~\ref{sec:res-est}). The tail is a property of ideal (instantaneous) edges: with a finite rise time
the coefficients decay faster (0.45\% outside $|h|\le12$ for a ramp of 0.3 slot), and truncated estimators recover
accordingly (Section~\ref{sec:res-est}).
\end{corollary}

\begin{corollary}[Energy on the original subcarrier]\label{cor:parseval}
For ideal $\pm1$ elements with a common pattern of duty $D$, $b^{[0]}=2D-1$ and
$\sum_h|b^{[h]}|^2=1$. A receiver that equalizes subcarrier $m$ with the diagonal of $\bm H$
alone (one-tap) collects the fraction $(2D-1)^2$ of the reflected energy and sees the rest
as interference; its signal-to-interference ratio is bounded by $(2D-1)^2/(1-(2D-1)^2)$, i.e.
$-4.8$\,dB at $D=0.75$ and $+1.1$\,dB at $D=0.875$. Joint equalization collects all of it,
an energy-collection gain of $(2D-1)^{-2}$: 6.0\,dB at $D=0.75$ (Appendix~\ref{app:parseval};
measured 5.9\,dB with the dispersive cell, Section~\ref{sec:res-valid}).
\end{corollary}

\begin{corollary}[Shared paths give low rank]\label{cor:lowrank}
Under A5, $a^{\rm d}_q(f)=\sum_{\ell=1}^{P_d}g_\ell e^{-j\bm k\bm u_\ell^{\top}\bm r_q}e^{-j2\pi f\tau_\ell}$,
so $[\bm a^{\rm d}_1,\dots,\bm a^{\rm d}_Q]$ has rank at most $P_d$, the number of user-hop paths,
whatever $Q$; the same holds for the input side. With a delay basis $\bm F\in\mathbb C^{M\times L}$
spanning delays $[0,\tau_{\max}]$, the element-domain factors can be written
$\bm a^{\rm d}_q=\bm F\bm U\bm\alpha_q$ and $\bm v_{q,s}=\bm F\bm V_s\bm\beta_q$ with
$\bm U,\bm V_s\in\mathbb C^{L\times r}$ and $r$ on the order of the number of paths.
\end{corollary}

\begin{corollary}[Groups, shifts and the direct path]\label{cor:groups}
(i) Elements with the same pattern share $\bm T_{q,s}$; with $G$ distinct patterns,
\begin{equation}
\bm H=\sum_{g=1}^{G}\sum_{s=1}^{2}\bm T_{g,s}\circ\big((\bm F\bm U)\bm\Lambda_g(\bm F\bm V_s)^{\top}\big),
\qquad\bm\Lambda_g=\textstyle\sum_{q\in\mathcal G_g}\bm\alpha_q\bm\beta_q^{\top}\in\mathbb C^{r\times r}.
\label{eq:model}
\end{equation}
(ii) If pattern $g$ is the base pattern cyclically shifted by $k$ slots,
$t_{g,s}[h]=t_{0,s}[h]e^{-j2\pi hk/K}$, i.e. $\bm T_{g,s}=\bm D_k\bm T_{0,s}\bm D_k^{*}$ with
$\bm D_k=\diag(e^{-j2\pi mk/K})$; this is the harmonic phase rule of space-time
coding~\cite{zhang2018}. (iii) Since each element is always in one of its two states,
$\bm T_{q,1}+\bm T_{q,2}=\bm I$; an unswitched (direct) path is the special case of equal
input responses in both states, so it is already contained in~\eqref{eq:model}.
\end{corollary}

\begin{proposition}[A clock offset is an opposite delay of the two hops]\label{prop:offset}
If the switching runs $d$ slots late relative to the receiver's timing, every kernel
becomes $\bm D\bm T_{g,s}\bm D^{*}$ with $\bm D=\diag(e^{-j2\pi md/K})$, and
\begin{equation}
(\bm D\bm T\bm D^{*})\circ(\bm a\bm b^{\top})=\bm T\circ\big((\bm D\bm a)(\bm D^{*}\bm b)^{\top}\big).
\label{eq:offset}
\end{equation}
Multiplying a frequency response by $\bm D$ delays it by $dT_u/K$; by $\bm D^{*}$, advances it.
\end{proposition}
\begin{proof}
Shifting the pattern by $d$ slots multiplies its $h$-th coefficient by $e^{-j2\pi hd/K}=
e^{-j2\pi\bar md/K}e^{+j2\pi md/K}$ with $h=\bar m-m$, which is~\eqref{eq:offset} entrywise.
\end{proof}

\begin{proposition}[Physical timing]\label{prop:physical}
Let the BS--RIS and RIS--user hops have bulk delays $\tau_b$ and $\tau_u$ on all their paths, and let the
receiver synchronize its FFT window to the first cascaded arrival. (i) If the RIS controller is locked to the
signal arriving at the RIS, the synchronized coupling matrix equals the one without bulk delays: bulk delays
have no effect. (ii) If the controller is locked to the BS clock, the synchronized coupling matrix equals the
one without bulk delays and with a clock offset $d=-\tau_b/T_{\rm slot}$.
\end{proposition}
\begin{proof}
The delay $\tau_b$ multiplies the input side by $e^{-j2\pi f_m\tau_b}$ and $\tau_u$ the output side by
$e^{-j2\pi f_{\bar m}\tau_u}$; synchronization multiplies the output by $e^{+j2\pi f_{\bar m}(\tau_b+\tau_u)}$. In (i) the
switching edges are delayed by $\tau_b$, which by Proposition~\ref{prop:offset} contributes
$e^{-j2\pi(\bar m-m)\Delta f\tau_b}$; all phases cancel. In (ii) the edges are not delayed, and the residual
phases are those of an offset $-\tau_b$. Both identities hold numerically to $7\cdot10^{-15}$
(Section~\ref{sec:res-timing}).
\end{proof}

\begin{remark}[When can the pilot residual identify the offset?]
By Proposition~\ref{prop:offset}, a receiver that assumes a wrong offset can still fit the pilots with a
delayed user hop and an advanced BS hop, \emph{if} its delay basis represents those shifted responses. With
the DPSS basis used here ($L=16$, $\tau_{\max}=0.3\,\mu$s) the projection error of an advanced response is
$1.1\cdot10^{-5}$, $2.7\cdot10^{-4}$, $1.4\cdot10^{-2}$ and $0.82$ for shifts of 50, 100, 200 and 500\,ns: the
residual is nearly flat only for small offsets, which are exactly the physical ones (Proposition~\ref{prop:physical}(ii)
gives $|d|T_{\rm slot}=\tau_b$, typically tens to hundreds of ns). A wrong offset then overfits the pilots:
it explains them with improbable responses and predicts held-out pilots poorly. Two criteria resist this
overfitting: the pilot residual of a full low-rank estimate evaluated on a fine grid, and the Bayesian evidence
of the linear first stage, which penalizes improbable responses explicitly. Section~\ref{sec:res-timing}
compares them; the evidence is the cheaper and smoother of the two, not the only one that works.
\end{remark}
